\hyphenation{iso-mor-phisms pa-pers ca-te-gory}

% Nice script font. My reference due to Ron Donagi and Paul Aspinwall

\DeclareFontFamily{U}{rsf}{}
\DeclareFontShape{U}{rsf}{m}{n}{
  <5> <6> rsfs5 <7> <8> <9> rsfs7 <10-> rsfs10}{}
\DeclareMathAlphabet{\mathscr}{U}{rsf}{m}{n}

\renewcommand{\mathcal}{\mathscr}

% Gothic font
%\DeclareMathAlphabet{\mathgth}{U}{euf}{m}{n}
%\usepackage{oldgerm}
%\DeclareMathAlphabet{\mathgth}{U}{ygoth}{m}{n}

%\DeclareFontFamily{U}{cyr}{}
%\DeclareFontShape{U}{cyr}{m}{n}{
%  <5> wncyr5 <6> wncyr6 <7> wncyr7 <8> wncyr8 <9> wncyr9 <10-> wncyr10}{}
%\DeclareMathAlphabet{\mathcyr}{U}{cyr}{m}{n}
%
%% Cyrillic font, for the TS group
%\input cyracc.def
%
%\newcommand{\Sha}{\mathcyr{\cyracc Sh}}

%%%%%%%%%%%%%%%%%%%%%%%%
% there should be a \makeatletter-free way of changing this, but I can't find
% it...

\makeatletter
\def\operator@font{\sf}
\makeatother
%\SetSymbolFont{operators}{normal}{OT1}{cmss}{m}{n}


%\setlength{\unitlength}{1 true in}

% 2 CATEGORIES%
%\newcommand{\gVar}{\mathgth{Var}}
%\newcommand{\gAlg}{\mathgth{Alg}}
%\newcommand{\gCat}{\mathgth{Cat}}
\newcommand{\gVar}{{\mathcal{V}\textsf{ar}}}
\newcommand{\gAlg}{{\mathcal{A}\textsf{lg}}}
\newcommand{\gCat}{{\mathcal{C}\textsf{at}}}
% Gothic letters
\newcommand{\gC}{\mathcal{C}}
% bold letters
\newcommand{\C}{\mathbf{C}}
\newcommand{\D}{\mathbf{D}}
% Script/Calligraphic letters
%\newcommand{\cA}{{\mathscr A}}
%\newcommand{\cB}{{\mathscr B}}
\newcommand{\cE}{{\mathcal E}}
\newcommand{\cF}{{\mathcal F}}
\newcommand{\cG}{{\mathcal G}}
\newcommand{\cH}{{\mathcal H}}
\newcommand{\cO}{{\mathcal O}}
% Random
\newcommand{\HKR}{{\mathsf{HKR}}}
\newcommand{\PBW}{{\mathsf{PBW}}}
\renewcommand{\k}{\textsf{k}}
% Math Operators
\DeclareMathOperator{\Tr}{Tr}
\DeclareMathOperator{\sTr}{Tr}
\DeclareMathOperator{\tr}{tr}
\DeclareMathOperator{\Nat}{Nat}
\DeclareMathOperator{\Hom}{Hom}
\DeclareMathOperator{\End}{End}
\DeclareMathOperator{\Td}{td}
\DeclareMathOperator{\ch}{ch}
\DeclareMathOperator{\ev}{ev}
\DeclareMathOperator{\coev}{coev}
%\DeclareMathOperator{\Id}{Id}
\DeclareMathOperator{\Ob}{obj}
\DeclareMathOperator{\Ext}{Ext}
\DeclareMathOperator{\sHom}{\underline{Hom}}
\DeclareMathOperator{\bHom}{\textbf{\textsf{Hom}}}
\DeclareMathOperator{\THom}{{2\mbox{-}Hom}}
\DeclareMathOperator{\OHom}{1\mbox{-}Hom}
\DeclareMathOperator{\HH}{HH}
\DeclareMathOperator{\K}{K}

\renewcommand{\H}{\mathsf H}
\newcommand{\id}{\mathsf {Id}}
\newcommand{\Id}{\mathsf {Id}}

% useful abbreviations


\newcommand{\iso}{\cong}
\newcommand{\dual}{\,{\stackrel{\chk}{\longleftrightarrow}}\,}

\newcommand{\chk}{{\scriptscriptstyle\vee}}
\newcommand{\op}{{\mathsf{op}}}
\newcommand{\cop}{{\mathsf{cop}}}
\newcommand{\pt}{{\mathsf{pt}}}
\newcommand{\bone}{{\mathbf 1}}
\newcommand{\adjoint}{\dashv}
\newcommand{\ra}{\rightarrow}
\newcommand{\Ra}{\Rightarrow}
\newcommand{\lra}{\longrightarrow}
\newcommand{\raiso}{\stackrel{\sim}{\longrightarrow}}
\newcommand{\scdot}{{\,\cdot\,}}
\newcommand{\MP}[2]{{\ensuremath{\left\langle #1,\, #2\right\rangle_{\rm M}}}}
\newcommand{\MPbig}[2]{{\ensuremath{\Biggl\langle #1,\, #2\Biggr\rangle_{\rm M}}}}
\newcommand{\SP}[2]{{\ensuremath{\left\langle #1,\, #2\right\rangle_{\rm S}}}}
\newcommand{\SPbig}[2]{{\ensuremath{\Biggl\langle #1,\, #2\Biggr\rangle_{\rm S}}}}
% This is used in order to have the heading in lowercase
%\renewcommand{\uppercase}{}
\renewcommand{\phi}{\varphi}
\newcommand{\SerF}{\mathsf{S}}
%\newcommand{\SerK}{\mathsf{s}}
\newcommand{\SerK}{\Sigma}
\newcommand{\kerto}{\to}
\newcommand{\open}{\textsf{O}}
\newcommand{\closed}{\textsf{C}}
\newcommand{\TwoCob}{2\textsf{Cob}_{\text{oc}}}
\newcommand{\TwoCobL}{\TwoCob^\Lambda}
\newcommand{\Tcirc}{\circ_2}
\newcommand{\circv}{\circ_{\text{v}}}
\newcommand{\circh}{\circ_{\text{h}}}
\newcommand{\one}{\mathbf{1}}
\newcommand{\OO}{\cO}
\newcommand{\blob}{\bullet}
% picture commands
\newcommand{\figdir}{pstex}
\newcommand{\pstex}[1]{\raisebox{-.45\height}{\input{\figdir/#1.pstex_t}}}
\newcommand{\pic}{\pstex}
\newcommand{\pichere}{\text{\fbox{pic}}}
% Other


\newarrow{TeXto}----{->}
\allowdisplaybreaks
\newcommand{\Desclabel}[1]{\textsf{#1:}}
\newenvironment{Description}{\begin{list}{}{\renewcommand{\makelabel}{\Desclabel}
               \setlength{\labelwidth}{2.5em}
               \setlength{\labelsep}{0.5em}
               \setlength{\itemindent}{2em}
               \setlength{\itemsep}{0.5em}% separation between items
               \setlength{\rightmargin}{0em}
               \setlength{\leftmargin}{1em}}}
               {\end{list}}


\newcommand{\pd}[1]{#1^{\dagger}}
\newcommand{\rote}{{\begin{turn}{270}$\!\!\!\epsilon\ $\end{turn}}}